\section{Hyperideal polyhedra: an actual local integrand}
\label{sec:hyperideal}

\subsection{From vertices to orthogeodesic lengths}

A strictly hyperideal vertex lies outside the closed ball in the projective
model, while every edge meets the ball. Its polar plane cuts off the vertex.
After all such truncations, the original faces form a bordered hyperbolic
surface; the truncation polygons give its geodesic boundaries. Ideal vertices
can be included by a limiting procedure and give cusps.

For a hyperideal tetrahedron choose unit spacelike representatives
$v_1,\ldots,v_4$ in Minkowski space of signature $(3,1)$. Choose their signs
so that they represent the vertices in one projective affine chart and
have negative pairwise inner products. Set
\begin{equation}
 q_{ij}=-\langle v_i,v_j\rangle=\cosh a_{ij}>1.
 \label{eq:qvertices}
\end{equation}
Here $a_{ij}$ is the length of the original edge between the two polar
truncation planes. It is an orthogeodesic of the intrinsic bordered surface.
The six $a_{ij}$ specify the vertex configuration up to ambient isometry on
the admissible tetrahedral domain.

Let $b_i^{jk}$ be the side of the truncation triangle at vertex $i$ that lies
in the original face $ijk$. The right-angled hexagon cosine rule gives
\begin{equation}
 \cosh b_i^{jk}=
 \frac{\cosh a_{ij}\cosh a_{ik}+\cosh a_{jk}}
      {\sinh a_{ij}\sinh a_{ik}},
 \qquad
 L_i=\sum_{\{j,k\}\subset\{1,2,3,4\}\setminus\{i\}}b_i^{jk}.
 \label{eq:truncationlengths}
\end{equation}
The four $L_i$ are intrinsic boundary lengths. They are not the sums of the
exterior dihedral angles around the hyperideal vertices.

These formulas also describe the domain concretely. The Gram matrix
\begin{equation}
 G_{ii}=1,\qquad G_{ij}=-q_{ij}\quad(i\ne j)
 \label{eq:gram}
\end{equation}
must have signature $(3,1)$. With all $q_{ij}>1$, the strict tetrahedral
condition is $\det G<0$. Indeed, each $3\times3$ principal determinant is
$1-q_{ij}^2-q_{ik}^2-q_{jk}^2-2q_{ij}q_{ik}q_{jk}<0$.
The trace and a negative $2\times2$ principal determinant then show that
this block has signature $(2,1)$. Interlacing and $\det G<0$ give
signature $(3,1)$ for $G$.

There is also a direct check of the geometric inequalities used below.
At vertex $i$, put
\[
 u_j=\frac{v_j+q_{ij}v_i}{\sqrt{q_{ij}^2-1}}\quad(j\ne i).
\]
These are unit timelike vectors in $v_i^\perp$, with
$-\langle u_j,u_k\rangle=\cosh b_i^{jk}>1$, so they lie in the same
sheet of its hyperboloid. Their Gram determinant is
$\det G/\prod_{j\ne i}(q_{ij}^2-1)<0$ by the Schur complement.
Thus they are three noncollinear points in $\HH^2$: their pairwise
distances $b_i^{jk}$ form a genuine triangle and satisfy strict triangle
inequalities. Conversely a Gram realization with this signature lies in
a common affine chart: a positive eigenvector for the largest eigenvalue
of $I-G$ gives a timelike positive combination of the $v_i$ whose inner
products with all four are negative. Every projective edge meets the ball
because $\langle v_i+v_j,v_i+v_j\rangle=2-2q_{ij}<0$.

It is possible to glue right-angled hexagons from a wider set of positive
arc-length assignments and obtain a valid intrinsic bordered surface. Such
arcs need not be the edges of its convex hyperideal realization. Thus the
polyhedral domain must be imposed when the purpose is to integrate a
canonical polyhedral cell.

\subsection{The WP density on a fixed-boundary leaf}

Mondello's Poisson formula for arc lengths has the boundary kernel
\begin{equation}
 K_L(s)=\frac{\sinh(L/2-s)}{\sinh(L/2)},\qquad 0<s<L.
 \label{eq:boundarykernel}
\end{equation}
For two different arcs $e,f$, the bracket in convention~\eqref{eq:fn} is
one half the sum of these kernels over their common boundary endpoints,
with oriented boundary distances from the endpoint of $e$ to that of $f$.
Changing the order changes the sign, since $K_L(L-s)=-K_L(s)$.
The factor $1/2$ results from evaluating the ordered wedge sum in
\cite[Theorem 4.1]{MondelloBoundary}.
That theorem is stated for a compact surface with strictly positive
geodesic boundaries and a maximal disjoint orthogeodesic arc system.
Boundary lengths are Casimirs of the Poisson tensor; inversion is performed
only after restricting to a fixed-boundary leaf. The matrix in all six
tetrahedral arc lengths is therefore not the matrix to invert.

Restrict now to the two-dimensional leaf where $L_1,\ldots,L_4$ are fixed.
The arcs $a_{12},a_{13}$ share just the boundary at vertex $1$, and the
relevant distance is $b=b_1^{23}$, up to orientation. Inverting the
two-dimensional Poisson tensor gives
\begin{equation}
 \boxed{
 \dV=
 2\frac{\sinh(L_1/2)}{\sinh(L_1/2-b_1^{23})}
 \abs{\dd a_{12}\wedge\dd a_{13}}.}
 \label{eq:hyperidealdensity}
\end{equation}
Since the three truncation sides form a genuine triangle,
$b_1^{23}<L_1/2$. Thus the denominator is positive, the bracket is nonzero,
and these two arc lengths are local coordinates on this leaf.

Formula~\eqref{eq:hyperidealdensity} is an actual expression for the volume
integrand in terms of the hyperideal geometry. Its remaining variables are
functions on the fixed-$L$ leaf; they cannot be held independently fixed
while $a_{12},a_{13}$ vary.

For an expression directly in the vertex Gram entries, put
\[
 D=q_{12}^2+q_{13}^2+q_{23}^2+2q_{12}q_{13}q_{23}-1.
\]
Equation~\eqref{eq:truncationlengths} gives
\[
 \cosh b=\frac{q_{12}q_{13}+q_{23}}
                   {\sqrt{q_{12}^2-1}\sqrt{q_{13}^2-1}},\qquad
 \sinh b=\frac{\sqrt D}
                   {\sqrt{q_{12}^2-1}\sqrt{q_{13}^2-1}}.
\]
Substituting in~\eqref{eq:hyperidealdensity} and using
$\dd q_{ij}=\sinh a_{ij}\,\dd a_{ij}$ yields
\begin{equation}
 \boxed{
 \dV=
 \frac{2\abs{\dd q_{12}\wedge\dd q_{13}}}
 {q_{12}q_{13}+q_{23}-\coth(L_1/2)\sqrt D}.}
 \label{eq:gramdensity}
\end{equation}
The positivity follows from the same truncation-triangle inequality.
For a concrete normalization check, take $q_{ij}=2$ for every edge.
Then $G=3I-2\mathbf1\mathbf1^t$ has determinant $-135$,
$a_{ij}=b_i^{jk}=a=\operatorname{arcosh}2$, and $L_i=3a$.
The kernel is $K_{3a}(a)=1/5$, so the bracket is $1/10$ and
\eqref{eq:hyperidealdensity} has coefficient $10$. In Gram coordinates
$D=27$ and $\coth(3a/2)=3\sqrt3/5$, giving coefficient $10/3$ in
\eqref{eq:gramdensity}, exactly as $\sinh^2a=3$ requires.
This evaluates a local density at one admissible point; it does not solve
the fixed-boundary equations throughout a cell.

The cusp limit requires renormalization. As a positive boundary shrinks,
the raw distances to its truncation plane diverge; subtract the endpoint
terms corresponding to a chosen limiting horocycle before comparing with
lambda lengths. If $s/L\to\rho$, the kernel itself has the regular limit
$K_L(s)\to1-2\rho$. A fixed-perimeter horocycle gauge is different from
the four-unit-lambda gauge in~\eqref{eq:lambdagauge}; changing that gauge
changes coordinate Jacobians, while preserving the descended WP form.
This is the cuspidal limiting procedure underlying
\cite[Theorem 4.5]{MondelloBoundary}.

\subsection{General cells and what remains to integrate}

The same procedure applies to any arc system. On a fixed-boundary leaf,
choose independent arc coordinates $x_1,\ldots,x_{2d}$ and form the
antisymmetric matrix
\[
 B_{ef}=\{x_e,x_f\}_{\WP}.
\]
If it is nonsingular, then
\begin{equation}
 \boxed{
 \dV=\frac{|\dd x_1\cdots\dd x_{2d}|}{|\Pf B|}
       =\frac{|\dd x_1\cdots\dd x_{2d}|}{\sqrt{\det B}}.}
 \label{eq:pfaffiandensity}
\end{equation}
This is simply the top exterior power of the inverse Poisson matrix. It
turns a local arc-bracket formula into a local volume formula in every
dimension.

Bao--Bonahon's angle characterization assumes a planar $3$-connected dual
graph, $0<\theta_e<\pi$, circuit sums at least $2\pi$ with equality
only around ideal vertices, and sums greater than $\pi$ on embedded paths
joining distinct vertices of one dual face without lying in its boundary
\cite[Theorems 1--2]{BaoBonahon}. The path condition must not be dropped
when passing from ideal to hyperideal cells.
For a complete genus-zero hyperbolic metric with finitely many cusps and
infinite-area ends, the global realization theorem gives a unique
finite-vertex convex generalized hyperbolic polyhedron, up to congruence.
This is the case without finite cone singularities of the precise
restatement in~\cite[Theorem 1.1]{FillastreMetrics}; Section 1.3 there
attributes its proof to Schlenker~\cite{Schlenker}. For a bordered surface,
one attaches funnels and then truncates along the polar planes.
Rivin's intrinsic theorem supplies
ideal existence and generalized uniqueness~\cite[Theorems 1.1--1.3]{RivinIntrinsic}.
The local Gram and Poisson calculations above are independent of invoking
this global existence theorem.

For fixed $L_i$, however, the angle cell is cut out further by the nonlinear
equations~\eqref{eq:truncationlengths}. The local density is therefore more
complicated than in the ideal case. The known total
\[
 V_{0,4}(L_1,L_2,L_3,L_4)=2\pi^2+\tfrac12\sum_iL_i^2
\]
provides a stringent target for a direct evaluation of these cell integrals.
No such full hyperideal evaluation is carried out in this monograph.

Budd's metric-map polynomials obey, for $g=0,1$,
\begin{equation}
 V_{g,n}^{(2\pi)}(\alpha_1,\ldots,\alpha_n)
 =2^{2-2g-n}V_{g,n}\bigl(\sqrt{\alpha_1^2-4\pi^2},\ldots,
                         \sqrt{\alpha_n^2-4\pi^2}\bigr).
 \label{eq:buddidentity}
\end{equation}
This does not extend to $g\ge2$. Nor does the displayed substitution describe
the actual relation between hyperideal angle perimeters and truncation
boundary lengths~\cite[Corollary 2 and Section 1.3]{Budd}. It is useful
evidence of a connection, with a separate geometric question behind it.

\subsection{A further tree-coordinate result}

There is now an additional concrete integration route for genus zero with
a distinguished cusp. Budd--Meeusen--Zonneveld construct a spine bijection
from a full-measure subset to explicitly decorated tree polytopes
\cite[Theorem 1]{BMZ}. Writing $m$ for the total number of ends, their
statement has the form
\begin{equation}
 \M_{0,m}^{\circ}(0,\lengths)\ \longleftrightarrow\
 \bigsqcup_{\mathfrak t}\mathcal A_{\mathfrak t}(\lengths),
 \qquad
 \dV=2^{m-3}\,\dd^{\,2m-6}u_{\mathfrak t}.
 \label{eq:treevolume}
\end{equation}
Thus its volume is a finite sum of ordinary polytope volumes, with this
specified constant density. Their construction uses spine and tree
coordinates rather than the Bao--Bonahon hyperideal angle cells. It
requires an origin cusp.
In~\eqref{eq:treevolume}, $u_{\mathfrak t}$ means the independent coordinates
used in~\cite[equation (5)]{BMZ}: two angles at each trivalent inner
vertex and all but the last entry in each of the two boundary simplices
of size $L_i/2$. It is not Euclidean hypersurface measure on the redundant
coordinate space. With $m=n+1$ in that paper, its factor $2^{n-2}$ is
precisely $2^{m-3}$ here.

Zonneveld's sequel removes that restriction using decompositions into
half-tight cylinders~\cite[Theorem 3]{Zonneveld}. For positive boundary
lengths, with a selected pair satisfying $L_1<L_2$, it gives a full-measure
bijection to tree and double-tree cells and identifies their WP measures.
The latter cells include a gluing length and a twist parameter. The resulting
finite integrations compute the general genus-zero volume polynomials;
equal-length volume values follow by continuity. These are further direct
combinatorial integrations, but they do not identify the spine cells with
the hyperideal dihedral-angle cells. The direct evaluation in those
polyhedral coordinates remains a distinct calculation proposed here.

